% !TeX program = lualatex
% Standalone notebook; compile twice: lualatex 30.tex
% Research batch 39 down to 25, 27 September 2026.
% Original section preserved verbatim except for marked insertion blocks.
\documentclass[11pt,a4paper]{article}
\usepackage[margin=24mm,headheight=14pt]{geometry}
\usepackage{fontspec}
\setmainfont{Latin Modern Roman}
\setsansfont{Latin Modern Sans}
\setmonofont{Latin Modern Mono}
\usepackage{amsmath,amssymb,mathtools}
\usepackage{unicode-math}
\setmathfont{Latin Modern Math}
\usepackage{microtype}
\usepackage{enumitem,xurl,needspace}
\usepackage[dvipsnames]{xcolor}
\usepackage{hyperref}
\hypersetup{unicode=true,colorlinks=true,linkcolor=MidnightBlue,urlcolor=MidnightBlue,
 pdftitle={Research notebook: Problem 30},
 pdfauthor={Open Problems Math research notebook},bookmarksnumbered=true}
\usepackage{bookmark,titlesec,fancyhdr}
\titleformat{\section}{\Large\bfseries\raggedright}{\thesection}{0.7em}{}
\pagestyle{fancy}\fancyhf{}
\fancyhead[L]{\small\sffamily Open Problems Math\enspace|\enspace Research notebook}
\fancyhead[R]{\small\sffamily Problem 30}
\fancyfoot[L]{\small\sffamily 27 September 2026}
\fancyfoot[R]{\small\sffamily \thepage}
\setlength{\parindent}{0pt}
\setlength{\parskip}{5pt plus 1pt}
\setlength{\emergencystretch}{3em}
\setcounter{secnumdepth}{1}
\setcounter{tocdepth}{1}
\setlist[itemize]{leftmargin=3em,itemsep=5pt,topsep=4pt,parsep=0pt}
\allowdisplaybreaks
\newcommand{\N}{\mathbb{N}}
\newcommand{\Z}{\mathbb{Z}}
\newcommand{\Q}{\mathbb{Q}}
\newcommand{\R}{\mathbb{R}}
\newcommand{\C}{\mathbb{C}}
\newcommand{\F}{\mathbb{F}}
\newcommand{\A}{\mathbb{A}}
\newcommand{\PP}{\mathbb P}
\newcommand{\E}{\mathbb{E}}
\newcommand{\eps}{\varepsilon}
\newcommand{\dd}{\,\mathrm d}
\newcommand{\sref}[2]{\hyperlink{source:#1:#2}{[S#2]}}
\newcommand{\eref}[2]{\hyperlink{extra:#1:#2}{[E#2]}}
\newif\ifshowcatalogue
\showcataloguetrue
\newcommand{\cataloguescope}[1]{\ifshowcatalogue
 \paragraph{Exact scope in the supplied catalogue.}
 \begin{quote}\small #1\end{quote}\fi}
\numberwithin{equation}{section}
\begin{document}
\setcounter{section}{29}
% BEGIN PRESERVED SECTION
% ================================================================
% problemId: problem.generalized-ramanujan-conjecture-for-gl-n
% source releaseRank: 30; source releaseStatus: open_with_solved_subcases
% exactTarget SHA-256: 5e9c74c4c5f4b9e9901374055377ac8acd3d1c9797b25948d5b652479d1f3794
\clearpage
\section{\texorpdfstring{Generalized Ramanujan Conjecture for GL\_n}{Generalized Ramanujan Conjecture for GL\_n}}\label{30}
\subsection{Definitions and mathematical statement}
Write a cuspidal automorphic representation as $\pi=\bigotimes'_v\pi_v$. A unitary irreducible local representation is tempered if it is weakly contained in the regular representation; for the groups in question this also has the usual matrix-coefficient characterization modulo the center. The conjecture is
\[
 \pi\text{ cuspidal on }\operatorname{GL}_n(\A_F),\quad
 \omega_\pi\text{ unitary}\quad\Longrightarrow\quad
 \pi_v\text{ tempered for every place }v.
\]
At an unramified finite place, in unitary normalization, this asks that all Satake eigenvalues have absolute value one. The target includes ramified and archimedean places as well.
\par\smallskip\textit{Formulation and concept references:} \sref{30}{1}.
\cataloguescope{Prove that for every number field F and every cuspidal automorphic representation pi of GL\_n(A\_F) with unitary central character, each local component pi\_v is tempered at every place v.}
\subsection{Short English statement}
Must every local component of a unitary cuspidal automorphic representation satisfy the strongest expected temperedness bound?
% BEGIN RESEARCH 30
\subsection{Research attempt}

\paragraph{Current frontier.}
The target concerns every number field, every cuspidal representation with
unitary central character, and every place, including ramified and
archimedean ones. Sarnak's formulation and distinction between local
unitarity bounds and temperedness are retained. \sref{30}{1}, Section~1.
Two substantial recent results apply to special arithmetic classes:
Boxer--Calegari--Gee--Newton--Thorne prove Ramanujan and Sato--Tate for
regular algebraic cuspidal $\operatorname{GL}_2$ representations of parallel
weight over CM fields, and Newton--Thorne establish all symmetric-power
liftings for regular-weight Hilbert modular forms. The former appeared in
2025 and the latter in the 2026 Annals volume. \eref{30}{1}, \eref{30}{2}.
These do not cover arbitrary nonalgebraic Maass representations or arbitrary
rank. The catalogue's $\operatorname{Sp}_4$ density reference concerns a
different group; no pointwise $\operatorname{GL}_n$ theorem is inferred from
it. Recent publication records and stated scopes were checked, without
claiming a new independent audit of the long automorphy proofs.

\paragraph{Attack route.}
At an unramified finite place, repeated tensor or symmetric powers magnify
any Satake eigenvalue of modulus greater than one. If appropriate powers
were automorphic and locally compatible, even a coarse uniform-in-rank
bound would rule out that eigenvalue. This is a classical motivation for
functoriality, not a newly discovered principle. \sref{30}{1}, Section~1.
The attempt below makes the precise amplification and cancellation tests
explicit. A second route uses all power traces of one local parameter;
its missing step is a subexponential-in-the-power estimate at that fixed
place. Neither route is allowed to discard ramified or infinite places
while claiming the full target.

\paragraph{Attempt.}
Fix $\pi$ and an unramified finite place $v$ of residue cardinality $q$.
In the source's unitary normalization let $\alpha_1,\ldots,\alpha_n\ne0$
be the Satake eigenvalues. Unitary central character gives
$\prod_i|\alpha_i|=1$. Set
\[
 a_r=\sum_{i=1}^n\alpha_i^r,\qquad R=\max_i|\alpha_i|.
\]

\textbf{Lemma 30.1 (cancellation does not hide the spectral radius forever).}
One has
\begin{equation}\label{r30:radius}
                  \limsup_{r\to\infty}|a_r|^{1/r}=R.
\end{equation}
\emph{Proof.} The upper bound follows from $|a_r|\le nR^r$.
Near zero the absolutely convergent generating function is
\[
 \sum_{r\ge1}a_r z^{r-1}=\sum_{i=1}^n\frac{\alpha_i}{1-\alpha_i z}.
\]
Group equal eigenvalues. A distinct eigenvalue $\alpha$ of multiplicity
$m$ contributes $m\alpha/(1-\alpha z)$, with nonzero residue $-m$ at
$z=\alpha^{-1}$. All other grouped terms are regular there, so none of
these poles cancels. The nearest pole to zero has modulus $1/R$.
The power-series radius of convergence, or directly the root test and
analyticity inside that radius, proves \eqref{r30:radius}. The shift from
$r-1$ to $r$ in the root limit does not change it.\hfill$\square$

It follows that the following fixed-place estimate would imply temperedness
at this $v$:
\begin{equation}\label{r30:powerbound}
 \text{for every }\varepsilon>0\text{ there is }C_{\pi,v,\varepsilon}
 \text{ such that }|a_r|\le C_{\pi,v,\varepsilon}q^{\varepsilon r}
 \quad(r\ge1).
\end{equation}
Indeed $R\le q^\varepsilon$ for every $\varepsilon>0$, hence $R\le1$;
the determinant condition then forces $|\alpha_i|=1$ for all $i$.
The constants may depend arbitrarily on $\pi,v,\varepsilon$, but not on
$r$. A bound only for one fixed $r$, or one whose constant secretly grows
exponentially in $r$, is insufficient.

There is also a direct link with the local logarithmic derivative:
\[
 L_v(s,\pi)=\prod_i(1-\alpha_iq^{-s})^{-1},\qquad
 -\frac{L_v'(s,\pi)}{L_v(s,\pi)}
       =(\log q)\sum_{r\ge1}a_rq^{-rs},
\]
where the expansion is initially justified by $q^{-\Re s}R<1$.
Extending this expansion into a larger half-plane without proof would
assume the bound being sought. Holomorphy of a completed global product
does not automatically make each individual Euler factor pole-free in
that region, because other factors can compensate.

\medskip
\textbf{Lemma 30.2 (a precise conditional amplification).}
Suppose an unbounded set of positive integers $m$ has the following
property: there exists an automorphic representation $\Pi_m$ on
$\operatorname{GL}_{\binom{n+m-1}{m}}$, isobaric as a sum of unitary
cuspidal representations, unramified at this particular $v$, whose
Satake multiset is exactly that of $\operatorname{Sym}^m(\alpha(\pi_v))$.
Then $\pi_v$ is tempered.

\emph{Proof.} For an unramified unitary generic cuspidal component, the
coarse estimate $|\beta|<q^{1/2}$ is available independently of its rank;
for a $\operatorname{GL}_1$ unitary character it is immediate. Sarnak
records this local bound in equation~(12), before the stronger
Luo--Rudnick--Sarnak estimate in equation~(14). \sref{30}{1}, pp.~7--8.
Taking the union of the Satake multisets of the isobaric summands gives
$|\beta|\le q^{1/2}$ for every eigenvalue of $\Pi_{m,v}$.
In the symmetric power, $\alpha_i^m$ is one of those eigenvalues.
Consequently $|\alpha_i|\le q^{1/(2m)}$ for an unbounded sequence of $m$.
Letting $m\to\infty$ and again using the determinant proves the claim.
No growth control on the degree of $\Pi_m$ is needed for this limiting
argument; the uniform exponent $1/2$ is what matters.\hfill$\square$

One must retain the unitary-isobaric hypothesis. General automorphic
representations may involve nonunitary twists in residual or induced
constructions, and the same estimate cannot simply be assigned to them.
Similarly, potential automorphy after a field extension has to be linked
back to the original local parameter; it is not identical to the
hypothesis as written. The cited recent symmetric-power results provide
such deep additional arithmetic structure in their own settings, not a
general proof of the hypothesis for arbitrary $\pi$. \eref{30}{1},
\eref{30}{2}.

\medskip
\textbf{Stress-testing low-moment and almost-everywhere substitutes.}
For any real $R>1$, the multiset
\[
               \{R,-R,R^{-1},-R^{-1}\}
\]
has determinant one and first power trace zero, yet is not tempered.
All its odd power traces vanish, while
$a_{2r}=2(R^{2r}+R^{-2r})$. Thus a cancellation argument based on the
first Hecke trace, or only odd powers, cannot establish
\eqref{r30:powerbound}. The companion script checks these identities at
rational $R$; the displayed formula already proves them for all $R>1$.
This is a model parameter, not a claim that an automorphic cuspidal
representation with this parameter exists.

The weak functoriality statement in catalogue Problem~7 matches local
parameters only outside a finite exceptional set $S_m$ that can depend
on the transfer. It does not imply that a fixed $v$ avoids $S_m$ for
unboundedly many $m$. For example, if finite places are enumerated as
$v_1,v_2,\ldots$, the sets $S_m=\{v_1,\ldots,v_m\}$ are all finite but
every fixed place belongs to all sufficiently large $S_m$.
Therefore even proving exactly that weak statement would not, by this
argument alone, establish the local compatibility in Lemma~30.2.
Additional strong transfer or another way to remove the exceptional
places is necessary. This is a quantifier obstruction, not a
counterexample to functoriality.

\paragraph{Stress test.}
The reduction above is deliberately restricted to unramified finite
places. There is no single Satake multiset in the same sense for an
arbitrary ramified component. Temperedness in Weil--Deligne or
Weil--$\operatorname{SL}_2$ language has to retain the monodromy data
and normalization; replacing it by unit modulus of every raw Frobenius
eigenvalue would be incorrect. At archimedean places the appropriate
condition is on real parts of Langlands exponents, not on a finite-place
Euler polynomial. The source explains these separate local conventions.
\sref{30}{1}, equations~(9)--(10) and the following paragraph.

Averaged estimates over primes do not bound the entire power sequence at
a fixed prime. Nor is the subexponential estimate supplied by the
unramified bound $|\alpha_i|\le q^{1/2-1/(n^2+1)}$: that gives a fixed
positive exponential rate in $r$, precisely the rate amplification is
trying to remove. The logarithmic-derivative manipulation is used only
in its absolutely convergent region. No unproved interchange or analytic
continuation of its Dirichlet series is inserted.

\Needspace{8\baselineskip}
\paragraph{Outcome.}
\textbf{Outcome: Exploratory approach with unresolved gap.}

The notebook proves an exact spectral-radius criterion and a conditional
symmetric-power amplification at each specified unramified place. It
identifies the fixed-place compatibility that weak almost-everywhere
functoriality does not provide, and gives an explicit low-moment
cancellation obstruction. These are deductions explaining a route and its
failure points; their historical novelty is not asserted. Neither the
required general family of locally compatible lifts nor the all-place
extension has been obtained. Thus the exact universal temperedness
statement, especially its ramified and archimedean components, remains
unresolved by this attempt.

% END RESEARCH 30

\subsection{Sources}
\begin{itemize}\raggedright
\item[\textnormal{[S1]}]\hypertarget{source:30:1}{} Peter Sarnak, Notes on the Generalized Ramanujan Conjectures, in Harmonic Analysis, the Trace Formula, and Shimura Varieties.
\url{https://web.math.princeton.edu/~sarnak/Preprints/SarnakFieldsNotes.pdf}.
{\small\textit{Catalogue formulation source.}}
\item[\textnormal{[S2]}]\hypertarget{source:30:2}{} A note on Sarnak's density hypothesis for Sp4.
\url{https://ems.press/content/serial-article-files/53056}.
{\small\textit{Catalogue research/status reference; not a proof endorsement.}}
\item[\textnormal{[S3]}]\hypertarget{source:30:3}{} Regularized spectral expansion of a Rankin-Selberg period and moments.
\url{https://arxiv.org/abs/2609.01358}.
{\small\textit{Catalogue research/status reference; not a proof endorsement.}}
% BEGIN ADDED REFERENCES 30
\item[\textnormal{[E1]}]\hypertarget{extra:30:1}{} George Boxer, Frank Calegari,
Toby Gee, James Newton and Jack A. Thorne, \emph{The Ramanujan and Sato--Tate
Conjectures for Bianchi modular forms}, Forum of Mathematics, Pi 13 (2025),
e10, published 21 February 2025. \url{https://doi.org/10.1017/fmp.2024.29}.
\url{https://www.repository.cam.ac.uk/items/45e39d5a-9515-4073-819e-8962395a7ece}.
{\small\textit{Abstract and theorem scope consulted 27 September 2026;
regular algebraic parallel-weight CM case, not arbitrary cuspidal forms.}}
\item[\textnormal{[E2]}]\hypertarget{extra:30:2}{} James Newton and Jack A.
Thorne, \emph{Symmetric power functoriality for Hilbert modular forms},
Annals of Mathematics 203 (2026), no.~1, 283--347.
\url{https://annals.math.princeton.edu/2026/203-1/p04}.
{\small\textit{Published scope and metadata consulted 27 September 2026.
Used for the regular-weight Hilbert modular setting only.}}

% END ADDED REFERENCES 30
\end{itemize}

% END PRESERVED SECTION
\end{document}
